% !TeX root = ../../Book5.tex
% 纲要标题：### C.3 \texorpdfstring{\(p\) 进群与连续表示}{p 进群与连续表示}
% 纲要位置：工作记录/计划纲要.md，第 4629 行。
% #### C.3.1 \texorpdfstring{\(p\) 进拓扑与光滑表示}{p 进拓扑与光滑表示}
% #### C.3.2 紧开子群的不变量与平均
% #### C.3.3 一个 Hecke 代数与非分歧特征标
% 节标题由计划纲要.md 中附录 C 的第 3 条要点整理。

\section{\texorpdfstring{\(p\) 进群与连续表示}{p 进群与连续表示}}
\label{b5:sec:C03}

% 纲要范围：
% * p 进群与连续表示的拓扑附加结构；群的算子代数参见第二卷附录 D

实李群中可以沿实参数微分；\(p\) 进群的拓扑则由开子群控制。研究它的复表示时，常见的“光滑”指每个向量被某个开子群固定，并不表示对实或复坐标求导。本节以紧开子群的平均和一个可完全计算的 Hecke 代数说明这一点。

\subsection{\texorpdfstring{\(p\) 进拓扑与光滑表示}{p 进拓扑与光滑表示}}
\label{b5:subsec:C03:01}

对素数 \(p\)，有理数上的绝对值规定为 \(|p^ru/v|_p=p^{-r}\)，其中 \(u,v\) 都不被 \(p\) 整除。其完备化为 \(\mathbb Q_p\)，单位球为 \(\mathbb Z_p\)。超度量不等式
\(|x+y|_p\le\max(|x|_p,|y|_p)\)
使 \(p^r\mathbb Z_p\) 同时开闭。展开 \(p\) 进数的各位数字可得
\[
\mathbb Z_p\cong\varprojlim_n\mathbb Z/p^n\mathbb Z.
\]
每个有限商紧，兼容数列在其乘积中构成闭集，故 \(\mathbb Z_p\) 紧。于是 \(\mathbb Q_p\) 是局部紧且完全不连通的拓扑域。

\ProseParagraphBreak
在 \(\operatorname{GL}_n(\mathbb Q_p)\) 中，\(\operatorname{GL}_n(\mathbb Z_p)\) 是紧开子群，而
\(1+p^rM_n(\mathbb Z_p)\)、\(r\ge1\)，组成单位元的开邻域基础。逆矩阵由收敛的几何级数给出，所以这些集合确实是子群。紧开子群可看成有限矩阵群的逆极限，故平均许多局部常值对象时仍能归结为有限商。

\begin{definition}\label{b5:def:smooth-admissible-padic}
设 \(G\) 为 Hausdorff 局部紧群，单位元处有紧开子群组成的邻域基础，\(V\) 为复向量空间，\(\pi:G\to\operatorname{GL}_{\mathbb C}(V)\) 为群表示。若每个 \(v\in V\) 的稳定子群开，就称 \((V,\pi)\) 为\term[guanghua-biaoshi]{光滑表示}[smooth representation]。若它光滑，并且每个紧开子群 \(K\subseteq G\) 的不变量空间 \(V^K\) 都有限维，就称为\term[kerongxu-biaoshi]{可容许表示}[admissible representation]。
\end{definition}

当 \(V\) 取离散拓扑时，光滑性等价于作用映射 \(G\times V\to V\) 连续。它不是任意 Hilbert 空间连续表示的同义词。有限维连续复表示限制到一个紧开投射有限子群时确实光滑。理由是一般线性群的邻域 \(\|A-I\|<1/2\) 不含非平凡子群：若一个矩阵的全部整数次幂都在其中，模不为一的特征值会使正幂或负幂离开邻域；模为一而不等于一的特征值也有某个幂离开；所有特征值为一时，非平凡 Jordan 块的幂无界。因此这样的矩阵只能为 \(I\)。连续性及开子群基础给出一个开子群，其像包含在这个邻域，故像平凡。这也表明无限维表示所需的拓扑不能省略。

\subsection{紧开子群的不变量与平均}
\label{b5:subsec:C03:02}

\begin{proposition}\label{b5:prop:compact-open-average}
设 \(K\) 为带逆极限拓扑的投射有限群，\(V\) 为复向量空间，\(\pi:K\to\operatorname{GL}_{\mathbb C}(V)\) 为光滑表示。对每个 \(v\in V\)，可选一个固定它的开正规子群 \(N\trianglelefteq K\)，并定义
\[
e_Kv=\frac1{[K:N]}\sum_{g\in K/N}\pi(g)v.
\]
此式与 \(N\) 及代表元选择无关，给出 \(V\to V^K\) 的投影。在光滑复表示范畴中，函子 \(V\mapsto V^K\) 正合。
\end{proposition}
\begin{proof}
一个向量的开稳定子群在紧群中有有限指数；取它有限多个共轭的交，得到开正规子群 \(N\)。代表元改变不影响 \(v\) 的像。若缩小到 \(N'\subseteq N\)，每个原陪集重复 \([N:N']\) 次，分母也同倍增大，故平均不变。两个选择都可与其交比较，所以良定义。对不同向量共同缩小正规子群，便可验证线性性。左乘置换有限商，故平均结果 \(K\) 不变；对已经不变的向量平均等于自身。

不变量函子保持核。对于满射 \(V\to W\)，若 \(w\in W^K\)，先任意取提升 \(v\)，则 \(e_Kv\) 是不变提升，因为平均与同态交换且 \(e_Kw=w\)。因此也保持满射，得到正合性。
\end{proof}

这正是有限群 Reynolds 平均在投射有限群上的逐向量版本，分母可以在 \(\mathbb C\) 中求逆。若改用特征 \(p\) 的系数，而有限商含 \(p\) 因子，平均分母不再可逆，正合性可能失败。

\begin{example}\label{b5:exm:smooth-not-admissible}
令加法群 \(G=\mathbb Q_p\) 在
\(V=\mathbb C[\mathbb Q_p/\mathbb Z_p]\)
上平移作用，方括号表示以陪集为基的有限支集向量空间。每个向量被开子群 \(\mathbb Z_p\) 固定，故表示光滑。但 \(\mathbb Q_p/\mathbb Z_p\) 无限，且 \(V^{\mathbb Z_p}=V\)，所以它不可容许。光滑性只控制每个向量的不变性邻域，不能控制整个固定空间的维数。
\end{example}

\subsection{一个 Hecke 代数与非分歧特征标}
\label{b5:subsec:C03:03}

取 \(G=\mathbb Q_p^\times\)、\(K=\mathbb Z_p^\times\)。每个元素唯一写成 \(p^nu\)，\(n\in\mathbb Z,u\in K\)，因此 \(G/K\cong\mathbb Z\)。令 \(\mathcal H(G,K)\) 为紧支集、双 \(K\) 不变的复值函数空间。把 \(K\) 的体积归一化为一，并以卷积
\[
(f_1*f_2)(g)=\int_G f_1(x)f_2(x^{-1}g)\,\mathrm{d}^\times x
\]
定义乘法。这里函数只在有限个 \(K\) 陪集上非零，积分可直接按有限陪集和计算，无须引入无限级数。它称为相对于 \(K\) 的\term[Hecke-daishu]{Hecke 代数}[Hecke algebra]。

\begin{proposition}\label{b5:prop:padic-torus-hecke}
设 \(p\) 为素数，\(G=\mathbb Q_p^\times\)、\(K=\mathbb Z_p^\times\) 带 \(p\) 进拓扑，\(\mathcal H(G,K)\) 为紧支集双 \(K\) 不变复值函数的卷积代数，所用乘法 Haar 测度使 \(K\) 的体积为一。对 \(n\in\mathbb Z\)，令 \(T_n\) 为 \(p^nK\) 的示性函数，\(T\) 为形式变量。则对 \(m,n\in\mathbb Z\) 有
\[
T_m*T_n=T_{m+n},\qquad
\mathcal H(\mathbb Q_p^\times,\mathbb Z_p^\times)
\cong\mathbb C[T,T^{-1}],\quad T_n\longmapsto T^n.
\]
对任意 \(a\in\mathbb C^\times\)，公式
\(\chi_a(p^nu)=a^n\)、\(n\in\mathbb Z\)、\(u\in K\)，给出 \(G\) 的一个光滑可容许一维复表示；\(T_n\) 在其 \(K\) 不变量上按 \(a^n\) 作用。
\end{proposition}
\begin{proof}
卷积被积函数非零，恰当 \(x\in p^mK\) 且 \(x^{-1}g\in p^nK\)，即 \(g\in p^{m+n}K\)。在这种情况下整个 \(p^mK\) 都贡献一，其乘法 Haar 体积为一；否则积分为零。所有双不变紧支集函数是有限个 \(T_n\) 的线性组合，且这些函数线性无关，故得到 Laurent 多项式代数。

乘法中赋值相加，故 \(\chi_a\) 是群同态；其核含开子群 \(K\)，所以光滑。一维空间对任何紧开子群的固定空间至多一维，故可容许。积分作用
\(\pi(T_n)v=\int_{p^nK}\chi_a(g)v\,\mathrm{d}^\times g=a^nv\)，给出最后结论。
\end{proof}

对 \(K\) 平凡的一维表示通常称为\term[feifenqi-tezhengbiao]{非分歧特征标}[unramified character]。这个算例把群表示中的参数 \(a\) 化成 Hecke 代数生成元 \(T\) 的特征值。更高秩群会出现非交换的卷积代数，以及双陪集和诱导表示的问题；此时不能把 \(G/K\) 一概当作群。Hilbert 空间表示和算子范数完备化需要第二卷附录 D 所述的群算子代数理论，所用拓扑条件与这里的光滑代数表示不同。
